\chapter{Frequent Urination}

\section*{Definition and Overview}
Frequent urination, medically called urinary frequency, is the need to pass urine much more often than usual. Most people urinate around 4--8 times a day, so a significant increase in that number, especially with small volumes each time or waking repeatedly at night, is a reason to take notice. The causes range from simple ones, such as drinking a lot of fluid or caffeine, to important medical conditions including urinary tract infections, diabetes, an enlarged prostate in men, overactive bladder, and pregnancy. The pattern of the frequency, the associated symptoms, and tests such as a urine check identify the cause. This chapter explains why urination becomes frequent and what to do about it.

\section*{Epidemiology}
Frequent urination is very common and affects people of all ages. It is a leading symptom of urinary tract infection, which affects around half of all women at some point, and of overactive bladder, which affects around 10--20% of adults and becomes more common with age. In men, an enlarged prostate affects the majority of men over 60 and commonly causes frequency and night-time urination. Diabetes, both undiagnosed and poorly controlled, causes frequency through high blood sugar. Frequent urination at night (nocturia) is particularly common in older adults, affecting more than half of those over 65.

\section*{Aetiology and Pathophysiology}
\begin{itemize}
    \item \textbf{High fluid intake:} drinking large volumes, especially of water, tea, coffee, or alcohol, increases urine production
    \item \textbf{Caffeine and alcohol:} act as diuretics, increasing urine output and bladder irritation
    \item \textbf{Urinary tract infection:} infection irritates the bladder, causing frequency, urgency, burning, and small volumes
    \item \textbf{Diabetes:} high blood sugar spills into the urine, drawing water with it, causing large volumes of urine and thirst
    \item \textbf{Overactive bladder:} the bladder muscle contracts involuntarily, causing frequency, urgency, and sometimes urgency incontinence
    \item \textbf{Enlarged prostate:} in men, the prostate compresses the urethra, making the bladder work harder and empty incompletely, causing frequency and nocturia
    \item \textbf{Bladder capacity problems:} a small bladder, or reduced capacity from surgery or radiotherapy
    \item \textbf{Medicines:} diuretics ("fluid tablets") increase urine output, as does some other medication
    \item \textbf{Pregnancy:} the growing uterus presses on the bladder, and hormones increase urine production
    \item \textbf{Neurological conditions:} stroke, Parkinson's disease, and spinal problems affect bladder control
    \item \textbf{Other causes:} pelvic organ prolapse, bladder stones, and rarely bladder cancer
\end{itemize}

\section*{Clinical Features}
\begin{itemize}
    \item Passing urine more than 8 times in 24 hours, with small or normal volumes
    \item Waking at night to pass urine (nocturia), once or repeatedly
    \item A sudden, strong urge to urinate (urgency)
    \item Burning, stinging, or pain on urination, suggesting infection
    \item Cloudy, strong-smelling, or blood-stained urine
    \item Excessive thirst and large volumes of urine, suggesting diabetes
    \item A weak stream, hesitancy, or dribbling in men with prostate problems
    \item Pelvic pressure or a feeling of incomplete emptying
    \item The pattern, onset, and associated symptoms point to the cause
\end{itemize}

\section*{Differential Diagnosis}
The causes span infections, metabolic, and bladder conditions.
\begin{itemize}
    \item Urinary tract infection, cystitis, and urethritis
    \item Diabetes mellitus and diabetes insipidus
    \item Overactive bladder and urgency incontinence
    \item Benign prostatic enlargement in men
    \item Pregnancy
    \item Excessive fluid, caffeine, and alcohol intake
    \item Medication effects, including diuretics
    \item Bladder stones, prolapse, and rarely bladder cancer
    \item Neurological and spinal conditions
\end{itemize}

\section*{When to Seek Medical Care}
Frequent urination that is new, persistent, painful, or accompanied by thirst, blood, or a change in stream should be assessed by a doctor. It is worth investigating because the underlying causes are treatable.

\textbf{Call 000 (or local emergency services) for:}
\begin{itemize}
    \item Inability to pass urine with severe lower abdominal pain (acute urinary retention)
    \item Blood in the urine with clots
    \item Fever, chills, and back or flank pain, which may indicate a kidney infection
    \item Sudden neurological symptoms with bladder problems
    \item Collapse or confusion with excessive thirst and urination
\end{itemize}

\section*{Diagnosis and Investigations}
\begin{itemize}
    \item \textbf{History:} fluid intake, medicines, pattern, and associated symptoms
    \item \textbf{Urine tests:} dipstick and microscopy for infection, blood, and glucose, with culture if infection is suspected
    \item \textbf{Blood tests:} blood sugar and HbA1c for diabetes, and kidney function
    \item \textbf{Bladder diary:} recording fluid intake and urine output over a few days
    \item \textbf{Flow studies and ultrasound:} measuring residual urine after voiding, especially in men with prostate symptoms
    \item \textbf{Further tests:} cystoscopy and urodynamic studies for selected cases
\end{itemize}

\section*{Management}
\begin{itemize}
    \item \textbf{Reduce fluid stimulants:} cut caffeine and alcohol, and moderate evening fluids, especially before bed
    \item \textbf{Treat infection:} antibiotics for urinary tract infection
    \item \textbf{Manage diabetes:} improve blood sugar control
    \item \textbf{Bladder retraining:} gradually lengthening the interval between voids and practising pelvic floor exercises
    \item \textbf{Overactive bladder medicines:} antimuscarinic and beta-3 agonist medicines reduce urgency and frequency
    \item \textbf{Prostate treatment:} medicines to relax the prostate or reduce its size, and surgery where needed
    \item \textbf{Adjust medicines:} review diuretics and time them so they least affect sleep
    \item \textbf{Pelvic floor physiotherapy:} for pelvic floor dysfunction and prolapse
    \item \textbf{Treat the specific cause:} bladder stones, neurological conditions, and other diagnoses have targeted treatments
\end{itemize}

\section*{Complications}
\begin{itemize}
    \item Sleep disruption and fatigue from nocturia
    \item Falls in older adults hurrying to the toilet at night
    \item Urgency incontinence when the bladder cannot hold on
    \item Progression of undiagnosed diabetes or kidney infection
    \item Reduced quality of life, work performance, and social participation
    \item Complications of urinary retention in severe prostate disease
\end{itemize}

\section*{Prognosis and Outcomes}
Most causes of frequent urination respond well to treatment. Urinary tract infections clear with antibiotics, overactive bladder improves with bladder retraining and medicines in the majority, prostate symptoms respond to medical and surgical treatment, and diabetes-related frequency settles with glucose control. Lifestyle changes, particularly reducing caffeine and evening fluids, help most people. With accurate diagnosis, the frequency is usually reduced and quality of life restored.

\section*{Prevention and Self-Care}
\begin{itemize}
    \item Moderate caffeine, alcohol, and large fluid volumes in the evening
    \item Keep a bladder diary to identify patterns before seeing a doctor
    \item Practise pelvic floor exercises regularly
    \item Empty the bladder fully and avoid rushing
    \item Drink enough fluid during the day, but time it sensibly
    \item Seek assessment for new, persistent, or painful frequency
    \item Review medicines and diabetes control with a doctor
    \item Treat constipation, which can worsen bladder symptoms
\end{itemize}

\section*{Sources and Further Reading}
The following reputable sources provide further information for healthcare students and patients seeking reliable, current advice about frequent urination.
\begin{itemize}
    \item Healthdirect Australia. \textit{Frequent urination.} https://www.healthdirect.gov.au/frequent-urination
    \item Continence Foundation of Australia. \textit{Frequent urination.} https://www.continence.org.au/
    \item Kidney Health Australia. \textit{Urinary symptoms.} https://kidney.org.au/
    \item NHS. \textit{Urinary frequency.} https://www.nhs.uk/
    \item Mayo Clinic. \textit{Frequent urination.} https://www.mayoclinic.org/
    \item MedlinePlus. \textit{Urination frequency.} https://medlineplus.gov/
\end{itemize}
